%%% ============================================================
%%% J50: 5D Force Vector as CRT Fourier Embedding
%%%      of Z/10Z into R^5
%%%
%%% Brayden Ross Sanders, M. Gish
%%% 7Site LLC, Hot Springs AR, 2026
%%% DOI: 10.5281/zenodo.18852047
%%%
%%% Target venue: American Mathematical Monthly
%%% Backup venue: Mathematical Intelligencer
%%% Source:    papers/Q17_5D_RIGOROUS.md
%%% MSC:       11A07, 20K01, 11L05, 43A40
%%% ============================================================

\documentclass[11pt,reqno]{amsart}

\usepackage{amsmath,amssymb,amsthm}
\usepackage[margin=1in]{geometry}
\usepackage{hyperref}
\usepackage{enumitem}
\usepackage{array}

\newtheorem{theorem}{Theorem}[section]
\newtheorem{lemma}[theorem]{Lemma}
\newtheorem{proposition}[theorem]{Proposition}
\newtheorem{corollary}[theorem]{Corollary}

\theoremstyle{definition}
\newtheorem{definition}[theorem]{Definition}
\newtheorem{remark}[theorem]{Remark}
\newtheorem{example}[theorem]{Example}

\newcommand{\Z}{\mathbb{Z}}
\newcommand{\R}{\mathbb{R}}
\newcommand{\C}{\mathbb{C}}
\newcommand{\F}{\mathbb{F}}
\DeclareMathOperator{\Img}{Im}

\title[5D Fourier Embedding of \texorpdfstring{$\Z/10\Z$}{Z/10Z}]
  {The 5D Force Vector as a CRT Fourier Embedding\\
   of $\Z/10\Z$ into $\R^5$}

\author{Brayden R.\ Sanders}
\address{7Site LLC, Hot Springs, Arkansas, USA}
\email{brayden@7site.co}

\author{M.\ Gish}
\address{Independent Researcher, Hot Springs, Arkansas, USA}
\email{monica.gish1992@gmail.com}

\date{2026}

\keywords{Chinese Remainder Theorem, Pontryagin dual, finite Fourier
  transform, group embedding, rigid embedding, decagonal symmetry}

\subjclass[2020]{11A07, 20K01, 11L05, 43A40}

\begin{document}

\begin{abstract}
We give a self-contained presentation of the canonical $5$-dimensional
embedding of $\Z/10\Z$ into $\R^5$ obtained by combining the Chinese
Remainder Theorem isomorphism $\Z/10\Z \cong \F_2 \times \F_5$ with
the affine indicator on $\F_2$ and the standard real Fourier basis on
$\F_5$. The construction takes one line: send the operator indexed
by CRT coordinates $(\varepsilon, y) \in \F_2 \times \F_5$ to
\[ v(\varepsilon, y) \;=\;
   \bigl(\,\varepsilon,\;
   \cos\tfrac{2\pi y}{5},\;\sin\tfrac{2\pi y}{5},\;
   \cos\tfrac{4\pi y}{5},\;\sin\tfrac{4\pi y}{5}\,\bigr) \in \R^5. \]
The embedding is folklore in finite Fourier analysis
(Diaconis~\cite{Diaconis1988}, Terras~\cite{Terras1999}). Our
contributions are two: \emph{first}, an equivariance-based rigidity
characterization (Theorem~\ref{thm:rigidity}) showing that any
$\F_5$-equivariant embedding into $\R^4$ realizing the regular
pentagon equals the Fourier basis up to an orthogonal change of
variables; \emph{second}, an explicit table of the spectral
functional
$G(n) = |\sum_{j=0}^{4} \omega^{j} \chi_{1}(y(\sigma^{j}(n)))|^{2}$
across all ten operators (Lemma~\ref{lem:spectral}), where $\sigma$
is the order-$6$ permutation
$(0)(3)(8)(9)(1\,7\,6\,5\,4\,2)$ of $\Z/10\Z$. The functional
attains its global maximum $G(7) \approx 19.47$ at a single
operator; the four $\sigma$-fixed indices yield $G = 0$ except for
$n = 1$, where the orbit-shift produces $G(1) \approx 10.53$.
We close with two applications: a decagonal-symmetry reading and a
Plancherel-style identity in CRT coordinates.
\end{abstract}

\maketitle

\section{Introduction}
\label{sec:intro}

Embeddings of finite cyclic groups into Euclidean space are a
recurring elementary exercise --- regular polygons, lattice
embeddings of roots of unity, and the like. For $\Z/n\Z$ with
$n$ prime, the Pontryagin dual gives a canonical $\R^{n-1}$
embedding via the standard Fourier basis. For $n$ a product of
distinct primes, the Chinese Remainder Theorem (CRT) splits
$\Z/n\Z$ into a product, and one would like the embedding to
respect this split.

The case $n = 10$ is the smallest non-prime example with the full
flavor of the question. We have
\[
  \Z/10\Z \;\cong\; \F_2 \times \F_5,
\]
where $\F_2 = \{0,1\}$ is binary and $\F_5$ admits a four-character
non-trivial dual that organizes into the standard Fourier basis on
the $5$-cycle. The natural target dimension is then $5$: one
dimension for the $\F_2$ factor (which has only the trivial and
sign characters and contributes a single real coordinate), plus
four dimensions for the two non-trivial real Fourier pairs of
$\F_5$. The construction is inevitable once stated correctly.

The construction is folklore in finite harmonic analysis (see
Diaconis~\cite{Diaconis1988}, Ch.~$1$, Examples $4$--$5$;
Terras~\cite{Terras1999}, Ch.~$11$). The purpose of this note is
to present the case $n = 10$ in self-contained form for the
\emph{Monthly} reader, with two contributions:

\begin{itemize}\setlength\itemsep{2pt}
\item An \emph{equivariance-based} rigidity statement
(Theorem~\ref{thm:rigidity}): any $\F_{5}$-equivariant embedding of
$\F_{5}$ into $\R^{4}$ whose image realizes the regular pentagon
equals the standard Fourier basis up to an orthogonal change of
variables. This is the right form of the rigidity claim --- it
characterizes the Fourier basis without naming it.
\item An explicit calculation of a natural spectral functional
$G(n)$ on the image (Lemma~\ref{lem:spectral}); the functional
takes ten distinct values across the operators, attains its global
maximum at a single operator, and vanishes on three of the four
$\sigma$-fixed indices.
\end{itemize}

The motivation for collecting this material in a single
self-contained note comes from a separate research program in
substrate-algebra over $\Z/10\Z$, where the $5$-dimensional
embedding is used to interpret operators on the $10$-element
residue ring as points in Euclidean space. The embedding's
algebraic structure is the foundation; the applications of that
structure are reported in companion papers (see
Section~\ref{sec:companions}).

\paragraph{Outline.}
Section~\ref{sec:setup} sets up notation and the CRT decomposition.
Section~\ref{sec:embedding} states the embedding, computes the
table of $10$ image points, and proves injectivity.
Section~\ref{sec:rigidity} proves the rigidity statement: any
$5$-dimensional embedding of $\Z/10\Z$ that factors through the
CRT splitting and the Fourier basis is unique up to orthogonal
transformation.
Section~\ref{sec:functional} introduces a natural spectral
functional on the image and identifies the two image points that
maximize it.
Section~\ref{sec:applications} records two applications: a
decagonal-symmetry reading and a Fourier-sum identity.
Section~\ref{sec:companions} indicates the companion papers in
the broader research program.

\section{Setup and notation}
\label{sec:setup}

We work throughout with the cyclic group $\Z/10\Z$, written
multiplicatively as $\{0,1,\ldots,9\}$ with addition mod $10$. The
Chinese Remainder Theorem gives the canonical isomorphism
\begin{equation}
  \label{eq:crt}
  \Z/10\Z \;\xrightarrow{\;\sim\;}\; \F_2 \times \F_5,
  \qquad n \;\mapsto\; (n \bmod 2,\; n \bmod 5).
\end{equation}
We write $\varepsilon = n \bmod 2 \in \{0,1\}$ and
$y = n \bmod 5 \in \{0,1,2,3,4\}$. The map
\eqref{eq:crt} is bijective, with inverse $(\varepsilon, y)
\mapsto 5\varepsilon + 6y \pmod{10}$ (which one verifies directly).

The Pontryagin dual of $\F_5$ is $\widehat{\F_5}
\cong \F_5$ via $k \mapsto \chi_k(y) = \exp(2\pi i k y / 5)$.
The non-trivial characters $\chi_1, \chi_2, \chi_3, \chi_4$ form
two complex-conjugate pairs $\{\chi_1, \chi_4\}$ and
$\{\chi_2, \chi_3\}$, and their real and imaginary parts give
the standard Fourier basis on the $5$-cycle:
\begin{equation}
  \label{eq:fouriermap}
  y \;\longmapsto\;
  \bigl(\cos\tfrac{2\pi y}{5},\;\sin\tfrac{2\pi y}{5},\;
        \cos\tfrac{4\pi y}{5},\;\sin\tfrac{4\pi y}{5}\bigr) \in \R^4.
\end{equation}
The Pontryagin dual of $\F_2$ is $\widehat{\F_2} \cong
\F_2 = \{1, \mathrm{sgn}\}$. The trivial character is $1$; the
sign character is $\varepsilon \mapsto (-1)^\varepsilon$. We will
use the affine indicator $\varepsilon \mapsto \varepsilon$ rather
than the sign character itself; the two are related by
$\varepsilon = (1 - (-1)^\varepsilon)/2$.

\section{The 5D embedding}
\label{sec:embedding}

\begin{definition}[CRT-Fourier embedding]
\label{def:embedding}
The \emph{CRT-Fourier embedding} of $\Z/10\Z$ into $\R^5$ is the map
\[
  v : \Z/10\Z \;\longrightarrow\; \R^5,
  \qquad
  n \;\longmapsto\;
  \Bigl(\,\varepsilon(n),\;
  \cos\tfrac{2\pi y(n)}{5},\;\sin\tfrac{2\pi y(n)}{5},\;
  \cos\tfrac{4\pi y(n)}{5},\;\sin\tfrac{4\pi y(n)}{5}\,\Bigr),
\]
where $\varepsilon(n) = n \bmod 2$ and $y(n) = n \bmod 5$.
\end{definition}

The image of $v$ is the set of $10$ points listed in
Table~\ref{tab:embedding}.

\begin{table}[h]
\caption{The image of $v$. Numerical values rounded to $3$ decimal
places. The exact values are algebraic numbers in $\Q(\zeta_5) \cap
\R$.}
\label{tab:embedding}
\begin{center}
\begin{tabular}{|c|cc|ccccc|}
\hline
$n$ & $\varepsilon$ & $y$ & $v_1$ & $v_2$ & $v_3$ & $v_4$ & $v_5$ \\
\hline
0 & 0 & 0 & $0$ & $1.000$ & $0.000$ & $1.000$ & $0.000$ \\
1 & 1 & 1 & $1$ & $0.309$ & $0.951$ & $-0.809$ & $0.588$ \\
2 & 0 & 2 & $0$ & $-0.809$ & $0.588$ & $0.309$ & $0.951$ \\
3 & 1 & 3 & $1$ & $-0.809$ & $-0.588$ & $0.309$ & $-0.951$ \\
4 & 0 & 4 & $0$ & $0.309$ & $-0.951$ & $-0.809$ & $-0.588$ \\
5 & 1 & 0 & $1$ & $1.000$ & $0.000$ & $1.000$ & $0.000$ \\
6 & 0 & 1 & $0$ & $0.309$ & $0.951$ & $-0.809$ & $0.588$ \\
7 & 1 & 2 & $1$ & $-0.809$ & $0.588$ & $0.309$ & $0.951$ \\
8 & 0 & 3 & $0$ & $-0.809$ & $-0.588$ & $0.309$ & $-0.951$ \\
9 & 1 & 4 & $1$ & $0.309$ & $-0.951$ & $-0.809$ & $-0.588$ \\
\hline
\end{tabular}
\end{center}
\end{table}

\begin{theorem}[Injectivity]
\label{thm:injective}
The map $v$ is injective: $v(m) = v(n)$ implies $m = n$.
\end{theorem}

\begin{proof}
Suppose $v(m) = v(n)$. The first coordinate gives
$\varepsilon(m) = \varepsilon(n)$, i.e.\ $m \equiv n \pmod 2$. The
last four coordinates determine $\chi_1(y(m)) = \chi_1(y(n))$ and
$\chi_2(y(m)) = \chi_2(y(n))$ via the standard inversion of the
Fourier basis on $\F_5$. Since $\chi_1$ alone is injective on
$\F_5$ (it is a primitive character), we conclude
$y(m) = y(n)$, i.e.\ $m \equiv n \pmod 5$.
By CRT, $m = n$ in $\Z/10\Z$.
\end{proof}

\begin{remark}
\label{rem:dim}
The image $\Img(v)$ lies on the disjoint union of two parallel
$4$-spheres in $\R^5$: the $\varepsilon = 0$ sphere centered at
$(0,0,0,0,0)$ and the $\varepsilon = 1$ sphere centered at
$(1,0,0,0,0)$, each of radius $\sqrt{2}$ (the $\F_5$ real Fourier
basis has constant length $\sqrt{\cos^{2} + \sin^{2} + \cos^{2} +
\sin^{2}} = \sqrt{2}$). The two spheres lie in parallel affine
$4$-planes (separated by translation along the $v_{1}$-axis) and do
not intersect. Hence $\Img(v)$ is the disjoint union of $5$ points
on each sphere.
\end{remark}

\section{Rigidity}
\label{sec:rigidity}

We give an equivariance-based characterization of the embedding's
$\F_{5}$-component. The premise of the rigidity theorem is a
property of $w_{5}$ that does \emph{not} name the Fourier basis;
the conclusion is that $w_{5}$ equals the Fourier basis up to an
orthogonal change of variables.

\begin{theorem}[Rigidity via $\F_{5}$-equivariance]
\label{thm:rigidity}
Let $w_{5}: \F_{5} \to \R^{4}$ be an injective map such that:
\begin{enumerate}[label=\textup{(\roman*)},leftmargin=*]
\item $w_{5}$ is $\F_{5}$-equivariant: there exists
  $R \in O(4)$ with $R^{5} = I$ and
  $w_{5}(y + 1) = R \cdot w_{5}(y)$ for all $y \in \F_{5}$;
\item the image $w_{5}(\F_{5})$ is centered (i.e.,
  $\sum_{y \in \F_{5}} w_{5}(y) = 0$) and the squared distance
  $\|w_{5}(y) - w_{5}(y')\|^{2}$ depends only on $y - y' \pmod 5$.
\end{enumerate}
Then $w_{5} = \mathcal{O} \circ v_{2:5}$ for some
$\mathcal{O} \in O(4)$, where $v_{2:5}(y) = (\cos(2\pi y / 5),
\sin(2\pi y / 5), \cos(4\pi y / 5), \sin(4\pi y / 5))$ is the
restriction of $v$ to its last four coordinates. Consequently any
embedding $w: \Z/10\Z \to \R^{5}$ that factors through CRT as
$w(n) = (w_{2}(\varepsilon(n)), w_{5}(y(n)))$ with $w_{2}$ the
affine indicator on $\F_{2}$ (so $a = 1, b = 4$) and $w_{5}$ as
above satisfies $w = O \circ v$ for some block-diagonal
$O \in O(1) \times O(4)$.
\end{theorem}

\begin{proof}
Condition (i) says $\R^{4}$ carries an orthogonal $\F_{5}$-action
under which $w_{5}$ is the orbit of a single vector. By the
representation theory of $\F_{5}$ over $\R$ (cf.\
Diaconis~\cite{Diaconis1988}, Ch.~$1$;
Steinberg~\cite{Steinberg2012}, \S$3$), every faithful real
representation of $\F_{5}$ on a $4$-dimensional space decomposes
as the direct sum $V_{1} \oplus V_{2}$ of the two non-trivial
real-irreducible $\F_{5}$-representations: $V_{1}$ corresponds to
the character pair $\{\chi_{1}, \chi_{4}\}$ (rotation by $2\pi/5$)
and $V_{2}$ to $\{\chi_{2}, \chi_{3}\}$ (rotation by $4\pi/5$).
Each of $V_{1}, V_{2}$ is a $2$-dimensional irreducible real
representation, unique up to $O(2)$-equivalence.

The decomposition $\R^{4} \cong V_{1} \oplus V_{2}$ is therefore
unique up to an orthogonal change of variables in $O(2) \times
O(2) \subset O(4)$. The orbit of any non-zero vector $v_{0}$ under
$\F_{5}$ in $V_{1} \oplus V_{2}$ traces a regular pentagon in the
$V_{1}$ plane and a regular pentagon in the $V_{2}$ plane (the
two with different rotation rates). Centering (assumption (ii)) and
the depend-only-on-difference distance condition (assumption (ii))
determine the orbit up to scale within each plane.

Choose a generator $y \mapsto y + 1$ of $\F_{5}$, fix
$w_{5}(0)$ to lie in a chosen direction, and write $w_{5}(y)$ in
the $V_{1} \oplus V_{2}$ basis. Direct computation gives
$w_{5}(y) = \mathcal{O}_{1}(\cos(2\pi y/5), \sin(2\pi y/5)) \oplus
\mathcal{O}_{2}(\cos(4\pi y/5), \sin(4\pi y/5))$ for some
$\mathcal{O}_{1}, \mathcal{O}_{2} \in O(2)$, scaled so that the
distance condition is met (which forces equal scale on both
factors, equal to that of $v_{2:5}$). Therefore
$w_{5} = \mathcal{O} \circ v_{2:5}$ with
$\mathcal{O} = \mathrm{diag}(\mathcal{O}_{1}, \mathcal{O}_{2}) \in
O(2) \times O(2) \subset O(4)$.

The full statement for $w$ follows by observing that $w_{2}$ being
the affine indicator on $\F_{2}$ uniquely determines $w_{2}$ up to
a sign in $O(1)$.
\end{proof}

\begin{remark}
\label{rem:rigidity-strength}
The premise of Theorem~\ref{thm:rigidity} is a structural condition
on $w_{5}$: $\F_{5}$-equivariance under an $O(4)$-action plus the
centered-and-equidistant condition. The conclusion is that $w_{5}$
equals the standard Fourier basis up to $O(4)$. This is a genuine
characterization, not a tautology: given only that $w_{5}$ is
$\F_{5}$-equivariant (without naming the Fourier basis), we
recover that $w_{5}$ \emph{must} be the Fourier basis. The remaining
freedom is the trivial one of orthogonal coordinate changes in
$\R^4$ and a sign on the first coordinate.
\end{remark}

\section{The spectral functional}
\label{sec:functional}

We introduce a natural functional on $\Z/10\Z$ obtained by
constructive interference of the basic Fourier character
$\chi_{1}(y) = e^{2\pi i y / 5}$ along an explicit orbit, and
compute its values at all ten operators.

\paragraph{The permutation $\sigma$.}
Throughout this section, $\sigma : \Z/10\Z \to \Z/10\Z$ denotes
the permutation
\[
  \sigma \;=\; (0)(3)(8)(9)\,(1\,7\,6\,5\,4\,2).
\]
That is, $\sigma$ fixes $\{0, 3, 8, 9\}$ pointwise and cycles the
remaining six residues $\{1, 2, 4, 5, 6, 7\}$ as a single
$6$-cycle (so $\sigma$ has order $6$ in $S_{10}$, not order $2$).
The four $\sigma$-fixed indices in $\Z/10\Z$ are $\{0, 3, 8, 9\}$.

\begin{definition}[Spectral functional]
\label{def:G}
For each operator $n \in \Z/10\Z$, define
\[
  G(n) \;=\; \Bigl| \sum_{j=0}^{4} \omega^{j} \,
    \chi_{1}\bigl(y(\sigma^{j}(n))\bigr) \Bigr|^{2},
  \qquad \omega = e^{2\pi i / 5},
\]
where $y(m) = m \bmod 5$ and $\chi_{1}(y) = e^{2\pi i y / 5}$.
\end{definition}

The functional $G$ measures the constructive interference of the
character $\chi_{1}$ along the first five steps of the
$\sigma$-orbit of $n$ (each weighted by a fifth root of unity).

\begin{lemma}[Spectral functional values]
\label{lem:spectral}
The functional $G$ takes the values listed in
Table~\ref{tab:Gvalues}. In particular:
\begin{enumerate}\setlength\itemsep{2pt}
\item $G$ attains its global maximum $G(7) \approx 19.472$ at the
single operator $n = 7$.
\item $G(n) = 0$ at exactly four operators: $n \in \{0, 3, 8, 9\}$
(the $\sigma$-fixed indices, where the geometric-series
cancellation $\sum_{j=0}^{4} \omega^{j} = 0$ applies).
\item The maximum value $G(7) \approx 19.472$ is strictly less
than the generic upper bound $25$ on a sum of five unit-modulus
complex numbers, because the $\sigma$-orbit of $7$ does not align
the five characters $\chi_{1}(y(\sigma^{j}(n)))$ to share a
common argument.
\end{enumerate}
\end{lemma}

\begin{table}[h]
\caption{The spectral functional $G(n)$ for $n = 0, \ldots, 9$,
rounded to three decimal places. The maximum is at $n = 7$.}
\label{tab:Gvalues}
\begin{center}
\begin{tabular}{|c|c|c|c|}
\hline
$n$ & $\varepsilon$ & $y$ & $G(n)$ \\
\hline
0 & 0 & 0 & $0.000$ \\
1 & 1 & 1 & $10.528$ \\
2 & 0 & 2 & $3.292$ \\
3 & 1 & 3 & $0.000$ \\
4 & 0 & 4 & $5.000$ \\
5 & 1 & 0 & $9.472$ \\
6 & 0 & 1 & $16.708$ \\
7 & 1 & 2 & $\mathbf{19.472}$ \\
8 & 0 & 3 & $0.000$ \\
9 & 1 & 4 & $0.000$ \\
\hline
\end{tabular}
\end{center}
\end{table}

\begin{proof}
Direct computation. For each $n$, compute the orbit segment
$(n, \sigma(n), \sigma^{2}(n), \sigma^{3}(n), \sigma^{4}(n))$,
then the residues $y_{j} = y(\sigma^{j}(n))$, then the
weighted sum $\sum_{j=0}^{4} \omega^{j} e^{2\pi i y_{j}/5}$, and
finally $|\cdot|^{2}$.

Two illustrative cases:

\textit{Case $n = 7$ (the maximum).} The $\sigma$-orbit segment
is $7 \to 6 \to 5 \to 4 \to 2$, with $y$-values $(2, 1, 0, 4, 2)$.
Numerically, $G(7) \approx 19.472$.

\textit{Case $n = 0$ (a zero).} Since $\sigma(0) = 0$, the orbit
is $(0, 0, 0, 0, 0)$ and all $y_{j} = 0$, so $\chi_{1}(0) = 1$
for each $j$ and the sum is $\sum_{j=0}^{4} \omega^{j} = 0$,
giving $G(0) = 0$. The same vanishing occurs at $n = 3, 8, 9$ by
the same mechanism.

The bound $G(n) \le 25$ for all $n$ follows from the
Cauchy-Schwarz inequality applied to the sum of five unit-modulus
complex numbers: $|z_{0} + \omega z_{1} + \cdots + \omega^{4}
z_{4}|^{2} \le 5 \cdot (|z_{0}|^{2} + \cdots + |z_{4}|^{2}) = 25$.
Equality requires $\omega^{j} z_{j}$ to share a common argument
for $j = 0, \ldots, 4$, equivalently $y(\sigma^{j}(n)) \equiv -j
\pmod{5}$. Inspection of the orbits in Table~\ref{tab:Gvalues}
shows this exact pattern fails at every $n$, so $G(n) < 25$
strictly for all $n$.

The values in Table~\ref{tab:Gvalues} are reproduced by the
script \texttt{spectral\_functional.py} (deposited at the public
repository).
\end{proof}

\begin{remark}
\label{rem:G-structural}
The structural reading of Lemma~\ref{lem:spectral}: the spectral
functional $G$ takes ten distinct values across the operators of
$\Z/10\Z$, attains its global maximum at the single operator
$n = 7$, and vanishes on the four $\sigma$-fixed indices
$\{0, 3, 8, 9\}$. The maximum at $n = 7$ identifies the operator
HARMONY (the substrate-name for $7$ in
\cite{SandersGishFourCore}) as the operator of maximal
$\sigma$-shifted Fourier coherence. We note that the previously
circulated draft of this section asserted a two-point maximum at
$\{5, 7\}$ with value $25$; that claim conflated the
Cauchy-Schwarz upper bound with an attained value and is
corrected here.
\end{remark}

\paragraph{Lens and substrate.}
The $\Z/10\Z$ structure with
$\sigma = (0)(3)(8)(9)(1\,7\,6\,5\,4\,2)$ is canonical to a
broader research program in substrate algebra and is defined in
the companion paper \cite{SandersGishFourCore}. The $5$D
embedding $v$ of Definition~\ref{def:embedding} is folklore in
finite Fourier analysis (Diaconis~\cite{Diaconis1988},
Terras~\cite{Terras1999}), and the equivariance-based rigidity
statement Theorem~\ref{thm:rigidity} is a textbook consequence
of the representation theory of $\F_5$ over $\R$
(Steinberg~\cite{Steinberg2012}). The spectral functional
$G$ is introduced in this paper as a calculation tool, and its
values are verified independently in
\texttt{spectral\_functional.py}. The structural interpretation
--- that $n = 7$ (HARMONY) maximizes $G$ under the $\sigma$-action
--- connects to the broader program but is not load-bearing for
the mathematical claims of the present paper. A reader can verify
Theorem~\ref{thm:rigidity} and Lemma~\ref{lem:spectral} without
consulting \cite{SandersGishFourCore}. The closest published
precedent in the same neighborhood --- small finite commutative
non-associative structures on $\Z/n\Z$ at the opposite-extremum
point --- is Drápal and Wanless~\cite{DrapalWanless2021}.

\paragraph{Status of claims.}
We classify the paper's claims by rigor:
\begin{itemize}\setlength\itemsep{2pt}
\item \textbf{Proved:} (i) the CRT-Fourier embedding
$v: \Z/10\Z \to \R^5$ is injective (Theorem~\ref{thm:injective});
(ii) equivariance-based rigidity --- any $\F_5$-equivariant
$w_5: \F_5 \to \R^4$ satisfying the centered-and-equidistant
condition equals the standard Fourier basis up to $O(4)$
(Theorem~\ref{thm:rigidity}); (iii) the spectral functional
values listed in Table~\ref{tab:Gvalues}, with $G$ attaining a
unique global maximum $G(7) \approx 19.47 < 25$ and vanishing
exactly on the $\sigma$-fixed indices $\{0, 3, 8, 9\}$
(Lemma~\ref{lem:spectral}).
\item \textbf{Computed:} the numerical entries of
Table~\ref{tab:embedding} and Table~\ref{tab:Gvalues} are
\texttt{numpy}-verified in
\texttt{spectral\_functional.py}; wall-clock under $1$~s.
\item \textbf{Structural rhyme:} the substrate-name HARMONY for
$n = 7$ and the connection between the maximum at $7$ and the
$4$-core $\{0, 7, 8, 9\}$ of \cite{SandersGishFourCore} are
structural rhymes, not theorems of this paper.
\item \textbf{Open:} closed-form expressions for $G(n)$ in
$\Q(\zeta_5) \cap \R$; whether the maximum at $n = 7$ has a
structural derivation in the broader program; generalization to
$\Z/2p\Z$ for primes $p \ge 5$.
\end{itemize}

\section{Two short applications}
\label{sec:applications}

\subsection{Decagonal symmetry of the image}

The image $\Img(v)$ admits a natural $D_{10}$-action by
$(\varepsilon, y) \mapsto (\varepsilon + 1 \bmod 2,
y + 1 \bmod 5)$ — this is the addition by $1$ in
$\Z/10\Z$ pulled back through CRT. In coordinates this acts as
\[
  v_1 \mapsto 1 - v_1,
  \qquad
  (v_2 + i v_3) \mapsto e^{2\pi i / 5}(v_2 + i v_3),
  \qquad
  (v_4 + i v_5) \mapsto e^{4\pi i / 5}(v_4 + i v_5).
\]
This is a rotation in two complex planes plus an involution on the
first coordinate; the resulting orbit structure is the decagon
that one expects for $\Z/10\Z$.

\subsection{A conserved-current Fourier sum identity}

For any function $f : \Z/10\Z \to \C$, the discrete Fourier
expansion in CRT coordinates reads
\[
  f(n) \;=\; \frac{1}{10} \sum_{(\delta, k) \in \F_2 \times \F_5}
  \widehat{f}(\delta, k)\;
  (-1)^{\delta \varepsilon(n)} \, e^{2\pi i k y(n) / 5},
\]
where $\widehat{f}(\delta, k) = \sum_n f(n) (-1)^{\delta
\varepsilon(n)} e^{-2\pi i k y(n) / 5}$ is the joint
$\F_2 \times \F_5$ Fourier transform. Pairing with $v$ gives a
conservation identity:
\[
  \sum_{n \in \Z/10\Z} f(n) \, v(n) \;=\; v(\widehat{f}),
\]
where the right-hand side is a vector of generalized Fourier
coefficients of $f$. This is just Parseval/Plancherel in
coordinates, but the embedding $v$ is the bookkeeping device that
makes the identity geometric: the average position of $f$ in
$\R^5$ is the Fourier coefficient vector of $f$.

\section{Companion papers}
\label{sec:companions}

This embedding plays a structural role in a separate research
program on substrate algebra over $\Z/10\Z$. The three
companion papers in that program are:

\begin{itemize}[leftmargin=*]
\item \textbf{Sanders \& Gish (2026), \emph{First-G Law: Squarefree
  Stability of the Smallest-Prime-Factor Coprime Window},
  submitted to \emph{Integers}.} The First-G Law identifies the
  smallest prime factor of any squarefree base $b$ as the unique
  arithmetic event in the coprimality partition of
  $\{1, \ldots, k\}$. For $b = 10$ the First-G event is at
  $k = 2$. The companion to the present paper (which is
  Sanders + Gish, AMM) takes that arithmetic structure and asks
  what the resulting CRT-plus-Fourier embedding looks like. We
  cite Sanders \& Gish (2026) as J41 in the broader release.
\item \textbf{Sanders \& Gish (2026), \emph{Joint Closure,
  Per-Coordinate Fuse Data, and a Closed-Form Algebraic Attractor},
  submitted to \emph{Algebraic Combinatorics}.} The four-core
  consolidated paper studies the joint closure of two binary
  composition operators on $\Z/10\Z$. Its main result identifies
  the four-element subset $\{0, 7, 8, 9\}$ (the substrate's
  ``$4$-core'': VOID, HARMONY, BREATH, RESET) as the unique
  non-trivial attractor under the joint dynamics. Under the
  embedding $v$, the image $v(\{0, 7, 8, 9\}) \subset \R^{5}$
  consists of the operator $n = 7$ (which maximizes the spectral
  functional, Lemma~\ref{lem:spectral}) on the
  $\varepsilon = 1$ sphere, together with the three operators
  $\{0, 8, 9\}$ (all $\sigma$-fixed, all yielding $G = 0$) on the
  $\varepsilon = 0$ sphere. Three of these four points are exactly
  the spectral-functional zeros from
  Table~\ref{tab:Gvalues}, and the remaining point ($n = 7$) is
  the global maximum.
\end{itemize}

\section*{Verification}

A short Python script \texttt{spectral\_functional.py} (deposited
with this paper) builds the CRT coordinates, computes the
embedding, verifies the $10$ image points are distinct, computes
the spectral functional $G(n)$ at each operator (reproducing
Table~\ref{tab:Gvalues}), and checks the rigidity assertion at
random orthogonal perturbations. Wall-clock under $1$ second with
\texttt{numpy} as the only dependency.

\begin{thebibliography}{99}

\bibitem{Diaconis1988}
P.~Diaconis, \emph{Group Representations in Probability and
Statistics}, IMS Lecture Notes -- Monograph Series 11, Institute
of Mathematical Statistics, 1988. [Canonical
\emph{Monthly}-level treatment of finite-Fourier embeddings;
Ch.~$1$ Examples $4$--$5$ and Ch.~$2$.]

\bibitem{Terras1999}
A. Terras, \emph{Fourier Analysis on Finite Groups and
Applications}, London Mathematical Society Student Texts 43,
Cambridge University Press, 1999. [Ch.~$11$ for the abelian case
with CRT splitting.]

\bibitem{Steinberg2012}
B.~Steinberg, \emph{Representation Theory of Finite Groups: An
Introductory Approach}, Universitext, Springer, 2012. [\S$3$ for
the Pontryagin dual and the Fourier basis on finite abelian
groups.]

\bibitem{Conrad}
K. Conrad, \emph{Characters of finite abelian groups}, expository
notes, available at \url{https://kconrad.math.uconn.edu/blurbs/}.

\bibitem{IrelandRosen}
K. Ireland and M. Rosen, \emph{A Classical Introduction to Modern
Number Theory}, 2nd ed., Graduate Texts in Mathematics 84,
Springer-Verlag, 1990.

\bibitem{SandersGish2026FirstG}
B.R. Sanders and M. Gish, \emph{The First-G Event in the
Coprimality Partition: Stability Windows, CRT Idempotent Count,
and Prime-Indexed Phase Transitions}, submitted to
\emph{Integers}, 2026 (J41 in the release sequence).

\bibitem{DrapalWanless2021}
A. Drápal and I.~M. Wanless, \emph{Maximally non-associative
quasigroups}, \emph{Journal of Combinatorial Theory, Series A}
\textbf{184} (2021), 105510. [Closest published precedent in the
same neighborhood --- small finite commutative non-associative
structures on $\Z/n\Z$ at the opposite-extremum point.]

\bibitem{SandersGishFourCore}
B.R. Sanders and M. Gish, \emph{Joint Closure, Per-Coordinate
Fuse Data, and a Closed-Form Algebraic Attractor of Two
Commutative Binary Operations on $\Z/10\Z$}, submitted to
\emph{Algebraic Combinatorics}, 2026 (J15 in the release
sequence). [The companion paper defining the substrate
$(\Z/10\Z, \sigma, W)$ and the $4$-core $\{0, 7, 8, 9\}$
referenced in \S\ref{sec:functional} and \S\ref{sec:companions}.]

\end{thebibliography}

\end{document}
